\section{Related Work}
\label{sec:related}

\subsection{Alignment Benchmark Design and Evaluation}

TruthfulQA \citep{lin2022truthfulqa} evaluates whether models produce truthful answers to questions that humans frequently answer incorrectly, measuring MC2 multiple-choice accuracy as a proxy for factuality. BBQ \citep{parrish2022bbq} is a question-answering benchmark specifically designed to detect social biases across nine demographic categories; it measures whether models give biased responses in ambiguous contexts. Both benchmarks represent alignment-relevant constructs --- truthfulness and bias avoidance --- that are distinct from general reasoning or knowledge recall.

The Llama~2 technical report \citep{touvron2023llama2} is among the most prominent examples of joint alignment benchmark reporting: TruthfulQA MC2, BBQ, and safety scores are reported for both base and chat variants, showing within-model improvement from RLHF fine-tuning. Our work extends this within-model perspective to the cross-model population: rather than comparing base vs.\ chat variants for a single model family, we characterize the pairwise correlation structure across $N{=}296$ diverse open-weight models.

MMLU \citep{hendrycks2021mmlu} measures broad academic knowledge across 57 tasks, serving as a standard proxy for general LLM capability. It has become a de facto scale indicator in leaderboard-based evaluations, and prior work has established its role as a general capability covariate ($R^2 = 0.32$ for AlpacaEval-LC rank variance in earlier pipeline iterations of this work).

\subsection{Benchmark Structure and Dimensionality}

Prior work has examined the latent structure of benchmark suites from a dimensionality perspective. BenchScope \citep{sha2026benchscope} characterizes the effective dimensionality (ED) of benchmark suites, reporting $\text{ED} \approx 1.7$ for the open LLM leaderboard --- suggesting roughly two latent evaluation axes. This finding is consistent with our result: if $\text{ED} > 1$, then at least some benchmarks measure constructs beyond a single capability axis, and alignment-specific co-movement can exist even after scale removal. BenchScope does not, however, provide a directional pairwise test of whether any specific covariate (such as MMLU) confounds a specific benchmark pair.

PCA-based analyses \citep{anon2026pca} decompose cross-model benchmark variance into principal components, finding that TruthfulQA loads primarily on PC2 (23.4\% of variance), a component orthogonal to PC1 (general capability). This is consistent with our partial correlation results: if TruthfulQA has PC2-type orthogonal variance, controlling for MMLU (PC1 proxy) should reduce raw correlations. Our approach uses partial regression rather than PCA, enabling a directed hypothesis test about MMLU's specific confound contribution to the TruthfulQA $\times$ BBQ pair.

\subsection{RLHF and Alignment Training Effects}

Reinforcement learning from human feedback (RLHF; \citealt{stiennon2020rlhf,ouyang2022instructgpt}) has become the primary post-training method for aligning language models with human preferences. InstructGPT \citep{ouyang2022instructgpt} demonstrates that RLHF substantially improves performance on human preference evaluations and helpfulness metrics. \citet{touvron2023llama2} document improvements on TruthfulQA and BBQ for Llama~2 Chat relative to the base model, suggesting that RLHF may jointly improve factuality and bias avoidance within a model family.

Prior iterations of this research pipeline established a confirmed RLHF improvement on TruthfulQA MC2 of +3.406 points (BCa CI [+2.589, +4.212]) across 321 base/chat pairs, consistent with the literature. Whether RLHF generalizes equally to bias avoidance benchmarks --- and thus whether RLHF fine-tuning drives cross-model alignment co-movement --- is a separate question. Our Tier~3 analysis addresses this directly; the null result (sign test $p = 0.686$ on bias proxy) suggests the cross-model residual partial correlation is unlikely to be primarily driven by RLHF effects, at least under our bias proxy.

\subsection{Scale Confounding in LLM Evaluation}

The problem of scale confounding in LLM evaluation has received increasing attention. Cross-model correlation studies that use raw Spearman or Pearson correlation are vulnerable to confounding by model scale: larger models tend to score higher on nearly all benchmarks, inflating apparent pairwise correlations. This is analogous to the classical multivariate confounding problem in causal inference, where a shared common cause (scale) induces spurious correlation between outcomes.

Partial correlation as a method for controlling confounders is well-established in biostatistics and social science \citep{liu2017spearman}, but has not been systematically applied to LLM benchmark evaluation. The \texttt{pingouin} library \citep{vallat2018pingouin} implements partial Spearman via the inverse covariance method, providing exact $p$-values and enabling cluster-bootstrapped confidence intervals --- the statistical approach we adopt. The Fisher $z$ transformation for testing differences between independent correlations \citep{fisher1915} provides the formal statistical test for whether MMLU control significantly changes a benchmark pair's correlation. To our knowledge, this combination --- partial Spearman + Fisher $z$ difference test --- has not been applied as an alignment benchmark diagnostic prior to our work.

\subsection{Our Position}

Our work sits at the intersection of these four strands. Unlike BenchScope and PCA-based dimensionality analyses, we provide a directional, hypothesis-testable diagnostic for a specific covariate's confound contribution to a specific benchmark pair. Unlike within-model RLHF studies, we characterize cross-model population-level correlation structure after scale removal. Unlike general partial correlation statistics, we explicitly apply the Fisher $z$ difference test as a gate for determining whether MMLU scale control is necessary before interpreting alignment benchmark co-movement.
